\begin{definition}
    Пусть $T \subset \R$. Случайная величина $\tau$ со значениями в $T \cup \{+\infty\}$ называется марковским моментом относительно фильтрации $\{\F_t\}_{t\in T}$, если имеет место включение $\{ \tau \leq t \} \in \F_t$ для всех $t$. Если $\tau < +\infty$ п.в., то $\tau$ называется моментом остановки.
\end{definition}

\begin{definition}
    Пусть $\{\F_t\}_{t \geq 0}$ -- фильтрация. Отображение $\tau \to [0, +\infty]$ называется опциональным моментом, если $\{\tau < t\} \in \F_t$
\end{definition}

\begin{example}
    Для процесса $\{\xi_n\}_{n \in \N}$ и порожденной им фильтрации для фиксированного борелевского множества $B \subset \R$ положим
    
    $\tau_B(\omega) = min\{n: \xi_n(\omega) \in B\}$, если такие $n$ есть, иначе $\tau_B(\omega) = +\infty$.
    
    Тогда $\tau_B$ — марковский момент, ибо для всякого $t \in \N$ множество $\{\tau_B \leq t\}$ равно объединению множеств $\{\omega: \xi_n(\omega) \in B\}, n = 1, \ldots, t$, входящих в $\F_{\leq t}$.

    Если $\F_t$ интерпретировать как класс событий, о наступлении или ненаступлении которых известно к моменту $t$, то марковость можно понимать так: если при каком-то $\omega$ марковский момент принял значение $t$, то к моменту $t$ это уже известно. В случае фильтрации, порожденной случайным процессом $\{\xi_t\}$, наглядный смысл таков: по наблюдениям значений процесса до момента $t$ включительно, можно сказать, наступил ли уже момент $\tau$. Например, попало ли $\tau$ в множество $B$ до момента времени $t$, можно судить по значениями процесса до момента $t$ включительно.
\end{example}

\begin{note}
    Следующая теорема об остановке говорит о неизменности матожидания мартингала в ограниченный момент остановки
\end{note}

\begin{theorem}
    Пусть $\{\xi_n\}_{n\in \Z_+}$ — мартингал относительно фильтрации $\{\F_n\}_{n \in \Z_+}$ и $\tau$ — ограниченный момент остановки. Тогда имеем $\E \xi_\tau = \E \xi_0$. В частности, если $\E \xi_0 = 0$, то $\E \xi_\tau = 0$.
\end{theorem}